% Compile with XeLaTeX or LuaLaTeX on an updated TeX Live 2024 or later.
\documentclass[11pt,a4paper]{article}
\usepackage[margin=25mm]{geometry}
\usepackage{latexbangla}
\setbanglamainfont{Kalpurush}
\setmainfont{TeX Gyre Termes} % Optional: change the English font.
\usepackage{amsmath,amsthm}
\usepackage[hidelinks]{hyperref}
\listfiles

% Document-specific environments belong in the document or its class.
\newtheorem{theorem}{উপপাদ্য}[section]

\begin{document}
\section{পিথাগোরাসের উপপাদ্য}
একটি আয়তাকার বাগানের দৈর্ঘ্য ৪ মিটার এবং প্রস্থ ৩ মিটার। বাগানের এক কোণ
থেকে বিপরীত কোণে সরাসরি যেতে কত দূর হাঁটতে হবে?\footnote{এখানে বাগানটিকে
সমতল ধরা হয়েছে এবং সব দৈর্ঘ্য একই এককে মাপা হয়েছে।}
উত্তর পেতে আমরা \textit{Pythagorean theorem} ব্যবহার করতে পারি।

The shortest route across the garden is its diagonal—বাংলায় যাকে বলি কর্ণ।
Imagine drawing that route on a plan: the two sides of the garden and the
diagonal form a right triangle. With sides of 3 and 4 metres, walking around
the edge takes 7 metres. But can we cross the garden in just 5 metres?
\textbf{A sketch suggests the answer; নিশ্চিত হতে চাই একটি প্রমাণ।}

\begin{theorem}
সমকোণী ত্রিভুজের দুই লম্ব বাহুর দৈর্ঘ্য $a$ ও $b$ এবং অতিভুজের দৈর্ঘ্য $c$ হলে,
\begin{equation}\label{eq:pythagoras}
  a^2+b^2=c^2.
\end{equation}
\end{theorem}
এই উপপাদ্যের একটি জ্যামিতিক প্রমাণ ইউক্লিডের \textit{Elements}-এ পাওয়া যায়~\cite{euclid}।

\subsection{বাগানের কর্ণ ও ক্ষেত্রফল}
\eqref{eq:pythagoras} থেকে বাগানের কর্ণের দৈর্ঘ্য পাই,
\begin{align*}
  c &= \sqrt{৩^{২}+৪^{২}}=\sqrt{২৫}=৫\text{ মিটার},\\
  \text{বাগানের ক্ষেত্রফল} &= \text{দৈর্ঘ্য}\times\text{প্রস্থ}
    =৪\times৩=১২\text{ বর্গমিটার।}
\end{align*}
\textbf{মনে রাখি: উপপাদ্যটি শুধু সমকোণী ত্রিভুজের জন্য প্রযোজ্য।}
বাগানের ধার ঘুরে হাঁটার বদলে কর্ণ বরাবর গেলে তাই ২ মিটার পথ কমে যায়।
\textit{A shorter route, the same destination—একই গন্তব্যে কম পথ।}

\autoref{tab:triples}-এ এমন আরও কয়েকটি বাহুর দৈর্ঘ্য দেওয়া হলো।
\begin{table}[ht]
  \centering
  \caption{সমকোণী ত্রিভুজের পূর্ণসংখ্যক বাহুর দৈর্ঘ্য}\label{tab:triples}
  \begin{tabular}{ccc}
    \hline
    প্রথম লম্ব বাহু & দ্বিতীয় লম্ব বাহু & অতিভুজ \\
    \hline
    ৩ & ৪ & ৫ \\
    ৫ & ১২ & ১৩ \\
    ৮ & ১৫ & ১৭ \\
    \hline
  \end{tabular}
\end{table}

\subsection{নিজে করি}
\begin{enumerate}
  \item দুই লম্ব বাহুর দৈর্ঘ্য ৬ ও ৮ একক হলে অতিভুজের দৈর্ঘ্য নির্ণয় করো।
  \item একটি আয়তক্ষেত্রের দৈর্ঘ্য ১২ এবং প্রস্থ ৫ একক।
    \begin{enumerate}
      \item আয়তক্ষেত্রটির কর্ণের দৈর্ঘ্য কত?
      \item কর্ণ বরাবর কাটলে প্রতিটি ত্রিভুজের ক্ষেত্রফল কত?
    \end{enumerate}
\end{enumerate}

\begin{thebibliography}{9}
  \bibitem{euclid} Euclid. \textit{Elements}, Book I,
  \href{https://mathcs.clarku.edu/~djoyce/elements/bookI/propI47.html}{Proposition 47}.
  Online edition by David E. Joyce.
\end{thebibliography}
\end{document}
